%======================================================================
\section{Compiling groups into games}\label{sec:compile}
%======================================================================

\Cref{sec:transfer} identifies the cost of a constraint game with the state-weighted profile of an algebra attached to
it. To give a game a prescribed cost, one needs a game whose algebra has the profile of a given group. This section
supplies the tools, all of which keep the state: algebraic maps between finitely presented algebras of unitaries
(\cref{sec:compile-algebra}), corner maps that recover a group inside a game algebra on the whole space
(\cref{sec:compile-corner}), a compiler of every finitely presented group with a central involution into a fixed boolean
constraint game (\cref{sec:compile-universal}), and three routes into linear system games (\cref{sec:compile-linear}).

\subsection{Presented algebras and algebraic maps}\label{sec:compile-algebra}

A \emph{presented algebra} $\mathcal A=\langle X\mid R\rangle$ is a finite set $X$ and a finite set $R$ of elements of the
group algebra $\C[F(X)]$ of the free group. Let $I$ be the two-sided ideal generated by $R\cup R^*$. The algebra is the
unital $*$-algebra $\C[F(X)]/I$, in which the generators are unitaries; it is an algebraic quotient, not a completion. A
\emph{model} of $\mathcal A$ at precision $e$ on $\C^d$ is a homomorphism $\varphi:F(X)\to\mathrm U(d)$ and a density
matrix $\rho$, with $\xi=\rho^{1/2}$, such that
\begin{equation}\label{eq:model}
\hs{P(\varphi)\xi}\le e\quad(P\in R),\qquad\hs{[\varphi(x),\xi]}\le e\quad(x\in X),
\end{equation}
where $P(\varphi)=\sum_uc_u\varphi(u)$ for $P=\sum_uc_uu$. Write $h_{\mathcal A}(e)$ for the least $d$, $h^S_{\mathcal A}(e)$
for the infimum of $S(\rho)$, and $\Lambda_{\mathcal A}(e)$ for the set of nonzero spectra of the density matrices of such
models, each spectrum a finite list of positive numbers.

\begin{example}\label{ex:presented}
(a) For a finite presentation $\langle S\mid R\rangle$ and a word $w$, the \emph{negative quotient}
$A^-=\langle S\mid\{r-1:r\in R\}\cup\{w+1\}\rangle$ is $\C[G]/\langle w+1\rangle$, and its models are those of
\cref{def:hlp-rho}; so $h_{A^-}(e)=\hlp^\rho(w;e)$. We write $h_G$, $h_G^S$ and $\Lambda_G$ for its profiles. For a solution
group, $w=J$, and $A^-_\Gamma$ is the negative quotient of~\cite{CM14,PS23}.

(b) For a boolean constraint system $B$ (\cref{def:bcs}), the presented algebra $\mathcal A(B)$ has a generator for each
variable and the relators $x^2-1$, $xy-yx$ for variables in a common context, and the ordered products $R_\alpha$ of
\cref{def:boolean-profile} for every forbidden assignment $\alpha$. It is the algebra of the game~\cite{PS23}. For a linear
system $(A,b)$ it is the negative quotient of the solution group, since in it a parity context forces its product to be
$\pm1$.
\end{example}

\begin{lemma}[Finite identities]\label{lem:finite-identity}
Let $(\varphi,\rho)$ be a model of $\langle X\mid R\rangle$ at precision $e$. Every $p=\sum_uc_uu\in\C[F(X)]$ has
$\opn{p(\varphi)}\le M_p$ and $\hs{[p(\varphi),\xi]}\le L_pe$, with $M_p=\sum_u|c_u|$ and $L_p=\sum_u|c_u||u|$. Every $P\in I$
has $\hs{P(\varphi)\xi}\le C_Pe$, where $C_P$ depends only on a chosen expression of $P$ as a finite sum of terms $aQb$ with
$a,b\in\C[F(X)]$ and $Q\in R\cup R^*$.
\end{lemma}

\begin{proof}
The first two bounds follow from the triangle inequality and \cref{lem:calculus}(a). For $Q\in R$,
$\hs{Q(\varphi)\xi}\le e$, and for $Q=P^*$ with $P\in R$, $\hs{Q(\varphi)\xi}=\hs{\xi P(\varphi)}\le\hs{P(\varphi)\xi}+
\hs{[P(\varphi),\xi]}\le(1+L_P)e$. For a term $aQb$ write $b(\varphi)\xi=\xi b(\varphi)+[b(\varphi),\xi]$; the term is at most
$M_a\bigl(M_b\hs{Q(\varphi)\xi}+M_QL_be\bigr)$. Sum over the terms.
\end{proof}

An element of $I$ is a polynomial for which such an expression exists; the expression is a finite certificate, and the
constants below depend on the chosen certificates, not on $d$ or $\rho$. For $N\ge1$ let $D_N$ replace each entry
$\lambda$ of a spectrum by $N$ entries $\lambda/N$, the spectrum of $I_N/N\otimes\rho$.

\begin{proposition}[Algebraic maps transfer profiles]\label{prop:algebraic-map}
Let $\mathcal A_1=\langle X_1\mid R_1\rangle$ and $\mathcal A_2=\langle X_2\mid R_2\rangle$ be presented algebras, $N\ge1$, and
let $A_x\in M_N(\C[F(X_1)])$, $x\in X_2$, be matrices such that $x\mapsto A_x$, $x^{-1}\mapsto A_x^*$ induces a unital
$*$-homomorphism $\mathcal A_2\to M_N(\mathcal A_1)$. There is a constant $C$ such that every model of $\mathcal A_1$ at precision
$e$ on $\C^d$ gives a model of $\mathcal A_2$ at precision $Ce$ on $\C^{Nd}$ with density matrix exactly $I_N/N\otimes\rho$.
Hence
\[
h_{\mathcal A_2}(Ce)\le N\,h_{\mathcal A_1}(e),\qquad h^S_{\mathcal A_2}(Ce)\le h^S_{\mathcal A_1}(e)+\log N,\qquad
D_N\bigl(\Lambda_{\mathcal A_1}(e)\bigr)\subseteq\Lambda_{\mathcal A_2}(Ce).
\]
\end{proposition}

\begin{proof}
Put $\xi_2=I_N/\sqrt N\otimes\xi$ and $\hat A_x=A_x(\varphi)$, evaluated entrywise. If every entry of $Y\in M_N(\C[F(X_1)])$
lies in $I_1$, then $\hs{Y(\varphi)\xi_2}^2=\frac1N\sum_{i,j}\hs{Y_{ij}(\varphi)\xi}^2=O(e^2)$ by \cref{lem:finite-identity};
likewise $\hs{[Y(\varphi),\xi_2]}=O(e)$ and $\opn{Y(\varphi)}=O(1)$ for every $Y$. Since the class of $A_x$ is unitary in
$M_N(\mathcal A_1)$, the entries of $A_x^*A_x-1$ and $A_xA_x^*-1$ lie in $I_1$, and \cref{lem:polar} gives unitaries $U_x$ with
$\hs{(U_x-\hat A_x)\xi_2}$, $\hs{\xi_2(U_x-\hat A_x)}$ and $\hs{[U_x,\xi_2]}$ all $O(e)$. Let $P\in R_2$, and let $P(U)$ and $P(\hat A)$
be its values with $x^{\pm1}$ replaced by $U_x^{\pm1}$, respectively by $\hat A_x$ and $\hat A_x^*$. The entries of the image
of $P$ lie in $I_1$, so $\hs{P(\hat A)\xi_2}=O(e)$. For each word $y_1\cdots y_m$ of $P$, with letters $Y_l=U^{\pm1}$ and
$Y'_l$ the corresponding $\hat A$ or $\hat A^*$,
\[
Y_1\cdots Y_m-Y'_1\cdots Y'_m=\sum_lY'_1\cdots Y'_{l-1}(Y_l-Y'_l)R_l,\qquad R_l=Y_{l+1}\cdots Y_m ,
\]
and $(Y_l-Y'_l)R_l\xi_2=(Y_l-Y'_l)\xi_2R_l+(Y_l-Y'_l)[R_l,\xi_2]$ with $R_l$ unitary. The first term is $O(e)$, using
$\hs{(U^*-\hat A^*)\xi_2}=\hs{\xi_2(U-\hat A)}$ for inverted letters, and so is the second. So
$\hs{P(U)\xi_2}=O(e)$ for every $P\in R_2$, and every $U_x$ has state commutator $O(e)$. The dimension is $Nd$, the state
is $I_N/N\otimes\rho$, and $S(I_N/N\otimes\rho)=S(\rho)+\log N$.
\end{proof}

The identity maps between two finite presentations of the same algebra satisfy the hypothesis with $N=1$, so profiles
depend on the presentation only through the constant in the precision. For instance a Tietze transformation of a group
presentation, or the passage between a solution group and the algebra of its linear system game, changes $e$ by a
constant factor and keeps the state.

\begin{lemma}[Rounding to involutions]\label{lem:boolean-models}
Let $B$ be a boolean constraint system with contexts of size at most $k$. A model of $\mathcal A(B)$ at precision $e$ gives,
by $B_x=\mathrm r(\varphi(x))$, involutions satisfying \ref{B1}--\ref{B3} of \cref{def:boolean-profile} with the same $\rho$ at
precision $C_ke$; conversely the involutions of \cref{def:boolean-profile} at precision $e$ form a model of $\mathcal A(B)$ at
precision $e$.
\end{lemma}

\begin{proof}
The converse is immediate, since $B_x^2=I$. For the first part, \cref{lem:rounding} gives $\hs{(B_x-\varphi(x))\xi}\le e$
and $\hs{[B_x,\xi]}\le3e$. Each relator of $\mathcal A(B)$ is a fixed polynomial of degree at most $k$ in contractions, and
replacing the letters $\varphi(x)$ by $B_x$ one at a time, as in \cref{lem:calculus}(e), costs $O(e)$ in the state.
\end{proof}

Polar rounding and the replacement of a presentation by another are standard in the theory of approximate
representations~\cite{SV18}. What \cref{prop:algebraic-map} adds is that both are done in one fixed state at linear
precision, that the images $A_x$ need be unitary only modulo the relations, which is why the two-sided estimate of
\cref{lem:polar} is needed, and that the whole spectrum is carried along.

\subsection{Corner maps}\label{sec:compile-corner}

The lower direction of a comparison needs a map from the group into the game algebra. Restricting a homomorphism keeps the
state but does not, by itself, send the distinguished involution to $-1$. The following maps do, by covering the
whole space with corners on which it is $-1$.

\begin{proposition}[Corner maps]\label{prop:corner}
Let $G=\langle S_G\mid R_G\rangle$ and $H=\langle S_H\mid R_H\rangle$ be finitely presented, with $w_G$ representing a central
involution of $G$ and $J$ one of $H$. Let $\psi_1,\dots,\psi_k:F(S_G)\to F(S_H)$ induce homomorphisms $G\to H$, put
$W_i=\psi_i(w_G)$, and suppose that, in $H$, the $W_i$ commute with each other and with every $\psi_j(s)$, $s\in S_G$,
$1\le i,j\le k$, and that some $t\in F(S_H)$ satisfies $tbt^{-1}=Jb$ for $b=W_1\cdots W_k$. In $A_H^-$ put
\[
P=\tfrac12(1-b),\qquad P_i=\tfrac12(1-W_i),\qquad Q_i=PP_i\prod_{j<i}(1-P_j),\qquad H_s=\sum_iQ_i\psi_i(s)Q_i ,
\]
and $F(s)=H_s+tH_st^{-1}$ for $s\in S_G$. Then $s\mapsto F(s)$ defines a unital $*$-homomorphism $A_G^-\to A_H^-$ with
$F(w_G)=-1$. Consequently there is a constant $C$ such that every model of $H$ at precision $e$ gives a model of $G$ at
precision $Ce$ on the same space with the same density matrix, and
\[
h_G(Ce)\le h_H(e),\qquad h^S_G(Ce)\le h^S_H(e),\qquad\Lambda_H(e)\subseteq\Lambda_G(Ce).
\]
For $k=1$ the hypothesis reads $[\psi_1(w_G),t]=J$, and $F(s)=P\psi_1(s)P+tP\psi_1(s)Pt^{-1}$.
\end{proposition}

\begin{proof}
Work in $A_H^-$, where $J=-1$. The $W_i$ are commuting self-adjoint unitaries, so the $P_i$ are commuting projections,
and $P$ and the $Q_i$ are projections. The $Q_i$ are mutually orthogonal, since for $i<j$ the factor $1-P_i$ of $Q_j$
annihilates $P_i$, and $\sum_iQ_i=P\bigl(1-\prod_i(1-P_i)\bigr)=P$, because $b\prod_i(1-P_i)=\prod_iW_i\frac{1+W_i}2=\prod_i(1-P_i)$
gives $P\prod_i(1-P_i)=0$. Every $Q_i$ is a polynomial in the $W_j$, so it commutes with every $\psi_l(s)$. Hence, for
words $u,v$ and with $H_u$ defined by the same formula, $H_uH_v=H_{uv}$: the cross terms contain $Q_iQ_j=0$. Also
$H_u^*=H_{u^{-1}}$, $H_r=P$ for $r\in R_G$ because $\psi_i(r)=1$ in $H$, and $H_{w_G}=\sum_iQ_iW_iQ_i=-P$ because
$W_iQ_i=(1-2P_i)Q_i=-Q_i$. Next $tPt^{-1}=\frac12(1-Jb)=1-P$, so $H_u=PH_uP$ and $tH_ut^{-1}=(1-P)tH_ut^{-1}(1-P)$, and the
cross terms of $F(u)F(v)$ vanish. So $F(u)F(v)=F(uv)$, $F(u)^*=F(u^{-1})$, $F(r)=P+(1-P)=1$ for $r\in R_G$, and
$F(w_G)=-1$. The images are finite polynomials, and \cref{prop:algebraic-map} with $N=1$ gives the rest. For $k=1$,
$Q_1=P$, and $[W,t]=WtW^{-1}t^{-1}=J$ is equivalent to $tWt^{-1}=JW$ for an involution $W$.
\end{proof}

For the flat profile, the lower bound of~\cite[Proposition~6.2]{SV18} uses a witness $t$ with $k=1$ to show that the image
of $w_G$ is far from $I$~\cite[Lemma~6.3]{SV18}, and then restricts the representation to the $-1$ eigenspace of that
image~\cite[Proposition~2.7]{SV18}; the restriction keeps a flat state flat but cuts out a sector. \Cref{prop:corner}
replaces the restriction by an algebraic map on the whole space: no sector is cut out, and no sector weight multiplies
the entropy. A group homomorphism cannot do the same here. If $\psi_1(w_G)$ were identified with $J$, the relation
$[\psi_1(w_G),t]=J$ would become $[J,t]=J$, which forces $J=1$ since $J$ is central. The case $k>1$ is used for the
extended presentations of \cref{prop:extended} and for the restricted linear compiler of \cref{sec:rational-linear}, where
several conjugates of the involution are needed to cover the space.

\subsection{A universal compiler into boolean constraint games}\label{sec:compile-universal}

The next theorem compiles every finitely presented group with a central involution into a fixed boolean constraint
game whose constraints are parity equations and one binary predicate, at a fixed cost in dimension and entropy.

\begin{theorem}[Universal compiler]\label{thm:universal}
Let $G=\langle S\mid R\rangle$ be a finite presentation with a generator $w\in S$ representing a central involution. There
is a boolean constraint system $B_G$, computable from the presentation, in which every context has at most three variables
and every constraint is either a parity equation or a binary constraint forbidding the assignment $(-1,-1)$, and there are
unital $*$-homomorphisms
\[
F:A_G^-\to\mathcal A(B_G),\qquad\sigma:\mathcal A(B_G)\to M_{12}(A_G^-),\qquad\sigma(F(a))=I_{12}\otimes a ,
\]
given by finitely many polynomials. Consequently there are constants $C,C'>0$ with
\[
h_G(Ce)\le h_{B_G}(e)\le12\,h_G(e/C'),\qquad h^S_G(Ce)\le h^S_{B_G}(e)\le h^S_G(e/C')+\log12 ,
\]
and $\Lambda_{B_G}(e)\subseteq\Lambda_G(Ce)$, $D_{12}(\Lambda_G(e/C'))\subseteq\Lambda_{B_G}(e)$. If $h_G(e)<\infty$ for every $e>0$,
the game of $B_G$, with a context chosen uniformly and then a variable of it uniformly, has quantum value one, and there are
$C_1,C_2>0$ with
\[
h_G(C_1\sqrt\eps)\le\Ent(G(B_G),\eps)\le12\,h_G(C_2\sqrt\eps)\qquad(\eps>0),
\]
and the same with $h^S_G$ and $\log12$ added on the right for the least entanglement entropy of pure strategies.
\end{theorem}

\emph{The first stage.} Let $B^0$ have involutive variables $e_1,e_2,e_3,A,T$ and $V_s$ for $s\in S$, and the following
conditions; write $P_i=\frac12(1-e_i)$ and $N=\{A\}\cup\{V_s:s\in S\}$.
\begin{enumerate}[label=(F\arabic*),itemsep=1pt]
\item\label{F1} One-hot frame: $e_1,e_2,e_3$ form a context whose allowed assignments are $(-,+,+)$, $(+,-,+)$, $(+,+,-)$. In
the algebra the $P_i$ are orthogonal projections with $P_1+P_2+P_3=1$.
\item\label{F2} Coordinate swaps: every $n\in N$ satisfies $ne_1n=e_2$, $ne_2n=e_1$, $ne_3n=e_3$, and
$Te_1T=e_3$, $Te_3T=e_1$, $Te_2T=e_2$.
\item\label{F3} Blank support: for every $n\in N$, $(e_3,n)$ is a context forbidding just $(-1,-1)$. In the algebra this
says $P_3(1-n)=0$, so $nP_3=P_3$.
\item\label{F4} Relators: put $U_s=V_sA$. For every $r\in R$, the word $r(U)$, in which $s$ is replaced by $V_sA$ and $s^{-1}$
by $AV_s$, conjugates $T$ to itself. This is imposed by a chain of conjugacy relations: if $r(U)=g_1\cdots g_n$ with letters
$g_l\in N$, new involutions $t_n=g_nTg_n$, $t_l=g_lt_{l+1}g_l$ for $1<l<n$, and the relation $g_1t_2g_1=T$.
\item\label{F5} Sign: $(V_w,A,e_3)$ is a context with $V_wAe_3=-1$.
\end{enumerate}
\emph{Upper map.} On three blocks over $A_G^-$ let $e_1,e_2,e_3$ be the diagonal sign matrices with $-1$ in block $1,2,3$,
let $A$ exchange blocks $1$ and $2$ and $T$ blocks $1$ and $3$, and put
\[
V_s=\begin{pmatrix}0&s&0\\s^*&0&0\\0&0&1\end{pmatrix},\qquad U_s=V_sA=\mathrm{diag}(s,s^*,1).
\]
For a word $r$, $r(U)=\mathrm{diag}(r(s),\tilde r(s),1)$, where $\tilde r$ replaces each letter by its inverse without
reversing the order. If $r=1$ in $G$, this commutes with $T$ \emph{whatever $\tilde r$ is}; the chain variables are mapped
to the corresponding conjugates of $T$, which are involutions. Finally $w=-1$ gives $U_w=-e_3$, which is \ref{F5}. All
conditions hold, so this is a unital $*$-homomorphism $\sigma_1:\mathcal A(B^0)\to M_3(A_G^-)$. The middle block need not
satisfy any relator, and it is not tested: the swap $T$ exchanges the first block with the third, which is fixed pointwise.

\emph{Lower map.} In $\mathcal A(B^0)$, $A$ and every $V_s$ exchange $P_1$ and $P_2$ and fix $P_3$, so $U_s$ commutes with the
$P_i$, and \ref{F3} gives $U_sP_3=P_3$; the same holds for $U_s^{-1}$ and every word $r(U)$. By \ref{F4} and $TP_3T=P_1$,
$r(U)P_1=Tr(U)P_3T=P_1$ for $r\in R$. So $\theta(s)=P_1U_sP_1$ is a unitary representation of $G$ in the corner with unit
$P_1$, and \ref{F5} gives $\theta(w)=-P_1$, since $e_3P_1=P_1$. Put
\[
F_1(s)=\theta(s)+A\theta(s)A+T\theta(s)T .
\]
The three terms live in the corners $P_1,P_2,P_3$, their cross products vanish and their units add up to $1$, so $F_1$
is a unital $*$-homomorphism $A_G^-\to\mathcal A(B^0)$ with $F_1(w)=-1$. Evaluated on $\sigma_1$ it gives
$\sigma_1(F_1(s))=I_3\otimes s$.

\emph{Second stage.} Apply the nice-conjugacy conversion and the six-equation gadget of Slofstra~\cite[Lemma~4.4 and
Proposition~4.2]{Slo19}, in the explicit form of~\cite[proof of Lemma~B.1, Steps~3 and~4]{DM-game}, to the conjugacy
relations of $B^0$, and keep the contexts \ref{F1}, \ref{F3} and \ref{F5} unchanged. The result is a boolean constraint system
$B_G$. The old variables generate a copy of $\mathcal A(B^0)$ in $\mathcal A(B_G)$, because the new equations imply the old
conjugacy relations by finite multiplication~\cite{Slo19,DM-game}, and the lifts of the two stages give a unital
$*$-homomorphism $\sigma_g:\mathcal A(B_G)\to M_4(\mathcal A(B^0))$ that sends every old variable $x$ to $I_4\otimes x$; it
respects the kept contexts because these involve only old variables. Finally replace the one-hot context by the parity
equation $e_1e_2e_3=-1$ and three binary constraints forbidding $(-,-)$ on the pairs: among commuting signs the parity
equation allows an odd number of negative signs and the binary constraints at most one, so the allowed assignments, and
the algebra, are unchanged.

\begin{proof}[Proof of \cref{thm:universal}]
Put $F=F_1$, followed by the inclusion of $\mathcal A(B^0)$, and $\sigma=(\mathrm{id}_4\otimes\sigma_1)\circ\sigma_g$. Then
$\sigma(F(s))=I_4\otimes I_3\otimes s$. Every context of $B_G$ has at most three variables: the gadget equations and \ref{F5}
have three, the binary constraints two. The profile bounds follow from \cref{prop:algebraic-map}, applied to $F$ with $N=1$
and to $\sigma$ with $N=12$, together with \cref{lem:boolean-models}. The game bounds follow from these and
\cref{thm:boolean-correspondence}; if $h_G(e)<\infty$ for every $e$, the return direction gives strategies whose losing
probability tends to $0$.
\end{proof}

Two clauses are essential. Without the one-hot frame, $e_1=e_2=e_3=-1$ and $A=T=V_s=1$ satisfy every other first-stage
condition and encode nothing. Without the blank support of one $V_s$, the source relators need hold only up to a sign
in the payload block: for $G=\langle u,v,q,w\mid u^3,\ uvq^{-1},\ w^2,\ [w,u],[w,v],[w,q]\rangle$, the assignment
$U_u=-I$, $U_v=U_q=I$, $U_w=\mathrm{diag}(-1,-1,1)$ in three blocks, with $V_s=U_sA$, satisfies every other first-stage
condition although $u^3$ and $uvq^{-1}$ are mapped to $-I$.

\begin{remark}[Sources and what is new]\label{rem:universal-sources}
The parts of this construction are known, with one exception.
\begin{itemize}[leftmargin=*,itemsep=1pt]
\item A map in the direction of $F$ exists in much greater generality: Slofstra's embedding theorem gives
$A_G^-\hookrightarrow A_\Gamma^-$ into a solution group~\cite[Theorem~3.1]{Slo20}, and Fanizza, Kroell, Mehta, Paddock, Rochette,
Slofstra and Zhao embed the group algebra of a group generated by involutions into the algebra of a boolean constraint
system, with relation decompositions of bounded size~\cite[Lemma~3.12]{FKMPRSZ25}; with~\cite[Proposition~3.13 and
Lemma~2.9]{FKMPRSZ25} this controls near-perfect strategies in non-tracial states. The relation-decomposition bound used there,
\cite[Lemma~2.9]{FKMPRSZ25}, is Lemma~3.9 of Mehta, Slofstra and Zhao~\cite{MSZ23}.
\item The second stage is Slofstra's~\cite[Lemma~4.4, Proposition~4.2 and Remark~4.5]{Slo19}; its lift of every old variable
to $I_4\otimes x$ is explicit in~\cite[Lemma~5.8 and Theorem~5.5]{FKMPRSZ25}, which also gives decomposition bounds and the
extension of traces. The chains of \ref{F4} are their flattened nested conjugacies~\cite[Definition~5.6 and Lemma~5.7]{FKMPRSZ25}.
\item The dilation $y\mapsto\mathrm{adiag}(y,y^*)$ behind $V_s$ is that of~\cite[Proposition~4.8]{Slo19}; the constraint language,
one-hot contexts and binary exclusions, is that of Paddock's encoding of synchronous games~\cite{Pad22}, recalled
in~\cite[Section~3.2]{PS23}; and the passage from one-hot to parity and exclusions is an instance of~\cite[Lemma~5.4]{PS23}.
\item For synchronous games, Harris's reduction to graph three-colouring~\cite[Lemma~4.4 and Theorems~4.5 and~4.6]{Har23} has a
whole-algebra split with no amplification: his six colour-permuted corners are orthogonal, central and add up to one. It
maps games to games.
\end{itemize}
What we have not found is the retraction $\sigma$ for an arbitrary finitely presented group: the blank third sector is what
lets the inverse sector of $V_s$ violate the source relators without being tested. It gives the upper bounds, and with them
the comparison of the whole spectrum in both directions.
\end{remark}

\subsection{Linear system games}\label{sec:compile-linear}

Three routes lead to linear system games. The first applies to virtual retracts of the presentations, with at most one
non-involutive generator, for which our companion paper writes out Slofstra's embedding~\cite[Appendix~B]{DM-game}, and keeps
the state in both directions; the second applies to Slofstra's extended presentations with any number of non-involutive
generators and keeps it in the lower direction; the third, the restricted compiler behind the explicit exported games, is
described in \cref{sec:rational-linear}.

\paragraph{Virtual retracts of $E\Gamma_0$ presentations.} A \emph{group over $\Z_2$} is presented as $\langle S:R\rangle$ with a
distinguished central involution $J\in S$ and $J^2,[J,s]\in R$. Let $\mathcal Q=\langle S_0\sqcup S_1:R_{\mathcal Q}\rangle$ be an
$E\Gamma_0$ presentation in the sense of~\cite[Definition~10.1]{DM-game}, following~\cite[Definition~4.7]{Slo19}, with group $H$,
and $w\in S_0$ central in $H$. We call $G$ a \emph{generator-preserving virtual retract of $(\mathcal Q,w)$ of index $k$} if there
are $H_0\le H$ of index $k$ and homomorphisms $\iota:G\to H_0$ and $\pi:H_0\to G$ with $\pi\circ\iota=\mathrm{id}$, such that each
$\iota(s)$ is a generator $x_s\in S_0\sqcup S_1$, with $x_s\ne x_{s'}$ for $s\ne s'$ and $x_J=w$. Suppose that $|S_1|\le1$ and that
every linear set of $\mathcal Q$ has at most three elements; reordering the product of a linear set, whose elements commute,
presents the same group, so we may also assume the ordering hypothesis of~\cite[Lemma~B.1]{DM-game}. That lemma, the explicit
form of Slofstra's embedding~\cite[Section~4]{Slo19}, gives a binary linear system $(A,b)$ with solution group $\Gamma$, a
homomorphism $\Psi$ from the free group on $S_0\sqcup S_1$ to $F(x_1,\dots,x_n,J)$ and a variable $x_{j_0}$ such that $\Psi(r)$
is trivial in $\Gamma$ for every $r\in R_{\mathcal Q}$ and $[\Psi(w),x_{j_0}]J$ is trivial in $\Gamma$.

\begin{theorem}\label{thm:compile}
Let $G$ be a generator-preserving virtual retract of index $k$ of $(\mathcal Q,w)$, where $\mathcal Q$ is an $E\Gamma_0$
presentation with $|S_1|\le1$ and linear sets of at most three elements, and let $\Gamma$ be the solution group
of~\cite[Lemma~B.1]{DM-game} for $\mathcal Q$ and $w$.
Write $h_G,h^S_G$ and $h_\Gamma,h^S_\Gamma$ for the profiles of $J$. There are constants $C,C'>0$ such that, for every $e>0$,
\[
h_G(Ce)\le h_\Gamma(e)\le8k\,h_G(e/C'),\qquad h^S_G(Ce)\le h^S_\Gamma(e)\le h^S_G(e/C')+\log(8k).
\]
If $h_G(e)<\infty$ for every $e>0$, the game $G(A,b)$ has quantum value one, and there are $C_1,C_2>0$ with
\[
h_G(C_1\sqrt\eps)\le\Ent(G(A,b),\eps)\le8k\,h_G(C_2\sqrt\eps)
\]
for every $\eps>0$; the same holds for the least entanglement entropy of pure strategies, with $h^S_G$ in place of $h_G$ and
$\log(8k)$ added on the right.
\end{theorem}

\begin{proof}
\emph{Upper bounds.} Let $t_1=e,\dots,t_k$ be coset representatives of $H_0$, and $c(h,i)=t_{h\cdot i}^{-1}ht_i\in H_0$ the
induction cocycle. For each generator $x$ of $\mathcal Q$ and each $i$ fix a word $\beta(x,i)$ for $\pi(c(x,i))$, with
$\beta(w,i)=J$; this is possible because $w=\iota(J)$ is central and lies in $H_0$, so that $w\cdot i=i$ and
$c(w,i)=t_i^{-1}wt_i=w$, whose image $\pi(w)=\pi(\iota(J))$ is $J$. The monomial
matrices $\sigma_0(x)$ with entry $\beta(x,i)$ in position $(x\cdot i,i)$ define a homomorphism from $H$ into the monomial
matrices over $G$, by the cocycle identity, and so a unital $*$-homomorphism $\sigma_0:A^-_{\mathcal Q}\to M_k(A_G^-)$ with
$\sigma_0(w)=-1$. The four steps of the proof of~\cite[Lemma~B.1]{DM-game} pass from $\mathcal Q$ through presentations
$G'$, $\widehat G$ and $\widehat G_1$ to $\Gamma$. Each sends every generator of the next presentation to a block-monomial
matrix with entries $\pm$ words in the previous generators, or constants; the first, third and fourth steps double the number
of blocks. The proof there shows that the image of every relator, apart from the relators $Jwu$ and $t^*wt^*=u$ of Step~2, is
block diagonal with block words that are trivial in the previous group~\cite[proof of Lemma~B.1, and Table~3]{DM-game}. Step~1 is a unital
$*$-homomorphism $\C[G']\to M_2(\C[H])$, but not a map of negative quotients: it sends each $x\in S_0$ to $\mathrm{diag}(x,1)$, so
$w$ goes to $\mathrm{diag}(-1,1)$ in $M_2(A^-_{\mathcal Q})$. Step~2 adds involutions $t^*,J,u$, sent to the matrix that exchanges
the two blocks, to $-1$ and to $-\mathrm{diag}(w,1)$; its relators $Jwu$ and $t^*wt^*=u$ hold because $w=-1$ in $A^-_{\mathcal Q}$.
So Steps~1 and~2 together give a unital $*$-homomorphism $A^-_{\widehat G}\to M_2(A^-_{\mathcal Q})$ with $J\mapsto-1$. Steps~3
and~4 send $J$ to $J\oplus J$ and give unital $*$-homomorphisms $A^-_{\widehat G_1}\to M_2(A^-_{\widehat G})$ and
$A^-_\Gamma\to M_2(A^-_{\widehat G_1})$. Their composite with $\sigma_0$ is a unital $*$-homomorphism $A_\Gamma^-\to M_{8k}(A_G^-)$,
given by finitely many words, and \cref{prop:algebraic-map} with $N=8k$ gives the upper bounds. Steps~3 and~4 are the lifts
of~\cite[Lemma~4.4 and Proposition~4.2]{Slo19}, which turn a $d$-dimensional approximate representation into a
$4d$-dimensional one restricting to four copies of it~\cite[Remark~4.5]{Slo19}; see also~\cite[Lemma~5.8 and
Theorem~5.5]{FKMPRSZ25}.

\emph{Lower bounds.} Let $\hat\iota(s)=x_s$ and $\psi=\Psi\circ\hat\iota$. It induces a homomorphism $G\to\Gamma$ that sends
$J$ to $\Psi(w)$, and $t=x_{j_0}$ satisfies $[\Psi(w),t]=J$ in $\Gamma$. As $J$ is central in $G$, \cref{prop:corner} with
$k=1$ gives the lower bounds.

\emph{The game.} If $h_G(e)<\infty$ for every $e$, so is $h_\Gamma(e)$, and \cref{thm:return} gives strategies whose losing
probability tends to $0$, so $\omega_q=1$. The two displayed game bounds follow from \cref{cor:correspondence} and the profile
bounds, and likewise for entropy.
\end{proof}

For the flat profile and $k=1$, the companion's embedding theorem gives the two directions with explicit
constants~\cite[Theorem~B.3]{DM-game}; \cref{thm:compile} adds virtual retracts, keeps the state in both directions, includes the
entropy, and works at the natural precision $\sqrt\eps$.

\begin{remark}[A complete instance]\label{rem:compile-instance}
For the group generated by involutions $a,b,c$ and an element $y$ with the relations $abc=1$, $[a,b]=[a,c]=[b,c]=1$,
$yay^{-1}=b$, $yby^{-1}=a$ and $ycy^{-1}=c$, and with distinguished involution $c$, which is a generator-preserving retract
of itself, the construction of~\cite[Lemma~B.1]{DM-game}, with the commutations inside the linear set kept as relators of
the solution group rather than repeated as conjugacy relations, gives a linear system with $223$ variables and $180$ equations,
as the counts of that lemma predict with no conjugacy relations among the involutions, and an $8$-block lift of every
variable. The script \texttt{verify/check\_compiler.py} builds it, checks all $942$ square, product, commutator and witness
relations of the lift exactly in the negative quotient of the source, and runs the strategy of \cref{thm:return} on
perturbed models. The group has exact finite-dimensional models, so this is a check of the construction, not a game with
growing entanglement.
\end{remark}

\paragraph{Extended presentations.} Slofstra's \emph{extended homogeneous-linear-plus-conjugacy presentations}~\cite[Section~6]{SV18}
(following~\cite{Slo19}) have generators $S=S_0\sqcup S_1$ with $S_1=\{y_1,\dots,y_l\}$ linearly ordered; the relations
$x^2=e$ for $x\in S_0$; for each set $V$ of a collection of ordered subsets of $S_0$, the relations $\prod_{x\in V}x=e$ and
$xy=yx$ for distinct $x,y\in V$; and otherwise only relations $xyx=z$ with $x,y,z\in S_0$, $xyx^{-1}=z$ with $x\in S_1$ and
$y,z\in S_0$, and $y_iy_jy_i^{-1}=y_j^{L_{ij}}$ with $j<i$ and $L_{ij}>0$. Let $w\in S_0$ represent a central involution of $G$
that is non-trivial in approximate representations. For such a presentation Slofstra and Vidick state that there are a
solution group $\Gamma$ and constants with $\hlp_\Gamma(J;2,\eps)\le N\,\hlp_G(w;2,C\eps)$~\cite[Proposition~6.2]{SV18}. Their proof
prints $N=8$: it takes from~\cite[Proposition~4.8]{Slo19} a $2d$-dimensional lift of a $d$-dimensional representation with
$w\mapsto(-I_d)\oplus I_d$, and then adds an involution $t$ and the central $J$ with $twt=wJ$, realized by the exchange of the two
blocks. That proposition eliminates the non-involutive generators one at a time, and each step doubles the dimension and lifts
every involutive generator $x$ as $\varphi(x)\oplus I$. For $l=1$ this is the lift used. For $l\ge2$ the marked involution becomes
$W=(-I_d)\oplus I_{(2^l-1)d}$, and no unitary $T$ satisfies $TWT^*=-W$, even approximately: $\Tr(TWT^*+W)=2\Tr W=2(2^l-2)d$, so
$\fn{TWT^*+W}\ge2(1-2^{1-l})$ by the Cauchy--Schwarz inequality, whatever the precision of $\varphi$. This concerns only the factor $N$ of the upper bound: the
lower bound of~\cite[Proposition~6.2]{SV18}, and the non-triviality of $J$ in approximate representations of $\Gamma$,
which~\cite{SV18} take from~\cite[Lemma~5.2]{Slo19}, do not use this lift. The lamp sources of \cref{sec:rational} have
$l=2$ or $l=2k$, and the group behind the Slofstra--Vidick game has $l=2$ (\cref{sec:transport-game}). We therefore use a variant
that changes the marked element before the sign extension, and prove what we use.

\emph{The balanced construction.} Let $l\ge1$ and $k=2^{l-1}$. Eliminating $y_1,\dots,y_l$ in order as
in~\cite[Proposition~4.8]{Slo19} adds involutions $f_j,v_j$ with $y_j=f_jv_j$ and gives a homogeneous-linear-plus-conjugacy
presentation $G'$, all of whose generators are involutions, with a homomorphism $G\to G'$ that is the identity on $S_0$. For
$0\le i<k$ let $g_i$ be the product of the $f_j$ over the binary digits $2^{j-1}$ of $i$, in decreasing order of $j$; put
$\psi_i(s)=g_isg_i^{-1}$ for $s\in S$, read in $G'$, and $W_i=\psi_i(w)$. Add to $G'$ the relations that $W_0,\dots,W_{k-1}$
commute with each other and with every $\psi_j(s)$, $s\in S$, $j<k$; an involution $p$ with the row $pW_0\cdots W_{k-1}=e$; the
central $J$; and an involution $u$ with the row $pu=J$ and the conjugacy $f_lpf_l=u$. Every $\psi_j(s)$ is a word in
involutions, so with auxiliary involutions for the conjugates, as in~\cite[Proposition~4.8]{Slo19}, all new relations are rows and
involutive conjugacies. The result $\widehat G$ is a linear-plus-conjugacy group, and~\cite[Lemma~4.4 and Proposition~4.2]{Slo19}
give a solution group $\Gamma=\Gamma(A,b)$ and a homomorphism $\widehat G\to\Gamma$ fixing $J$. In the lift below the $W_i$ are
$-1$ on disjoint blocks, and the new marker $p$ has its two eigenvalues with equal multiplicities.

\begin{proposition}\label{prop:extended}
Let $G$ and $w$ be as above, with $l\ge1$ non-involutive generators, let $\Gamma=\Gamma(A,b)$ be the solution group of the
balanced construction, and put $N_G=2^{l+2}$. There are constants $C_1,C_2>0$ such that, for the game $G(A,b)$ with either
question distribution of \cref{def:lsg} and every $\eps>0$:
\begin{enumerate}[label=(\alph*),itemsep=1pt]
\item every strategy that wins with probability at least $1-\eps$ has Schmidt rank at least $h_G(C_1\sqrt\eps)$, and the
nonzero spectrum of Bob's reduced density matrix lies in $\Lambda_G(C_1\sqrt\eps)$;
\item if $\hlp_G(w;2,e)<\infty$ for every $e>0$, then $\omega_q(G(A,b))=1$ and $\Ent(G(A,b),\eps)\le N_G\,\hlp_G(w;2,C_2\sqrt\eps)$.
\end{enumerate}
\end{proposition}

\begin{proof}
(a) In $\Gamma$ each $\psi_i$ induces a homomorphism $G\to\Gamma$, the $W_i=\psi_i(w)$, $i<k$, commute with each other and with
every $\psi_j(s)$, and $b=W_0\cdots W_{k-1}$ equals $p$, so $t=f_l$ satisfies $tbt^{-1}=u=Jb$. \Cref{thm:transfer} gives a model of
$\Gamma$ at precision $O(\sqrt\eps)$ on the Schmidt support, with Bob's reduced state, and \cref{prop:corner} with these $k$ maps
$\psi_i$ and this $t$ turns it into a model of $G$ with the same state.

(b) Let $\varphi$ be a $d$-dimensional $e$-representation of $G$ with $\varphi(w)=-I$, and let $\gamma$ be the
$2^ld$-dimensional $O(e)$-representation of $G'$ that the $l$ lifts of~\cite[proof of Proposition~4.8]{Slo19} give: at step
$j$, with $\rho$ the representation reached so far, every involutive generator $x$ of the current presentation becomes
$\rho(x)\oplus I$, $f_j$ exchanges the two summands,
$v_j=\bigl(\begin{smallmatrix}0&\rho(y_j)^*\\\rho(y_j)&0\end{smallmatrix}\bigr)$, and every later $y_i$ becomes $\rho(y_i)\oplus\rho(y_i)$.
Number the $2^l$ blocks of size $d$ so that step $j$ adds the blocks $2^{j-1},\dots,2^j-1$. By induction on the steps, $\gamma(f_j)$
exchanges the blocks $i$ and $i+2^{j-1}$ for $i<2^{j-1}$ and fixes the others; every generator in $S$, read in $G'$, acts
block-diagonally; and $\gamma(w)$ is $-I$ on block $0$ and $I$ on the others. Hence $\gamma(g_i)$ carries block $0$ to block $i$,
every $\gamma(\psi_j(s))$ is block-diagonal, and $\gamma(W_i)$ is $-I$ on block $i$ and $I$ elsewhere. Put
$\gamma(p)=\gamma(W_0)\cdots\gamma(W_{k-1})$, which is $-I$ on the first $k$ blocks and $I$ on the others, $\gamma(u)=-\gamma(p)$ and
$\gamma(J)=-I$, and give the auxiliary involutions their defining values. The added relations then hold exactly; in particular
$\gamma(f_l)$ exchanges the two halves, so $\gamma(f_l)\gamma(p)\gamma(f_l)=-\gamma(p)=\gamma(u)$. So $\gamma$ is an
$O(e)$-representation of $\widehat G$ with $J\mapsto-I$, and~\cite[Remark~4.5]{Slo19} turns it into a $2^{l+2}d$-dimensional
$O(e)$-representation of $\Gamma$ with $J\mapsto-I$. So $\Gamma$ has flat models with $J\mapsto-I$ at every precision, of dimension at
most $N_G\,\hlp_G(w;2,Ce)$, and \cref{thm:return} with the flat state turns them into strategies.
\end{proof}

For $l=1$ the construction gives $N_G=8$, the value printed in~\cite{SV18}. A fixed $N_G$ adds $\log N_G=l+2$ to the entropies of
the upper strategies and leaves every leading coefficient and exponent unchanged. The restricted compiler of \cref{sec:rational-linear}
has the same architecture with its own conversion stages; \cref{prop:extended} does not certify its specifications. For the group
of~\cite[Proposition~6.2]{SV18} itself the lower direction needs only the relation $[\Psi(w),z]=J$ in its solution group, which holds
there; \cref{sec:transport} uses it with the corner map for $k=1$. The upper bound for that group comes from a lift specific to its
models, proved in our companion paper~\cite[Lemma~11.13]{DM-game} (\cref{sec:transport-game}).
